\documentclass[12pt,a4paper]{article}
\usepackage{fontspec}
\setmainfont{Times New Roman}
\usepackage[top=2cm,bottom=2cm,left=2.5cm,right=2cm]{geometry}
\usepackage{enumitem}
\usepackage{tabularx}
\usepackage{booktabs}
\usepackage{array}
\newcolumntype{Y}{>{\centering\arraybackslash}X}
\pagestyle{plain}
\linespread{1.15}

\begin{document}
% --- HEADER 2 CỘT CHUẨN NGHỊ ĐỊNH 30/2020/NĐ-CP ---
\noindent
\begin{minipage}[t]{0.45\textwidth}
    \begin{center}
        \fontsize{11pt}{13pt}\selectfont
        CÔNG AN THÀNH PHỐ HÀ NỘI\\
        \textbf{PHÒNG CẢNH SÁT HÌNH SỰ (PC02)}\\[-0.2em]
        \rule{3.2cm}{0.6pt}\\[0.25cm]
        \fontsize{11pt}{13pt}\selectfont Số: \textbf{00/HS}
    \end{center}
\end{minipage}
\hfill
\begin{minipage}[t]{0.52\textwidth}
    \begin{center}
        \fontsize{11pt}{13pt}\selectfont
        \textbf{CỘNG HÒA XÃ HỘI CHỦ NGHĨA VIỆT NAM}\\
        \textbf{Độc lập - Tự do - Hạnh phúc}\\[-0.2em]
        \rule{3.6cm}{0.6pt}\\[0.25cm]
        \textit{Hà Nội, ngày 25 tháng 07 năm 2016}
    \end{center}
\end{minipage}

\vspace{0.6cm}

\begin{center}
    \textbf{\Large TÀI LIỆU VỤ ÁN}\\[0.15cm]
    \textit{}
\end{center}

\vspace{0.4cm}
\textit{"Hãy để tôi được chuộc lại tội lỗi của mình. Nếu không vì hành động đêm đó đã đẩy Khang đến nguy hiểm, thì đứa bé đã không phải rơi vào cảnh mồ côi. Mọi ân oán lẽ ra nên để lại trong quá khứ. Vì mọi sinh linh bé nhỏ đều là vô tội, 20 năm trước hay bây giờ cũng như vậy"}


Cầm bức thư trên tay, ánh mắt điên dại của Hà chợt khựng lại. Lần đầu tiên, Hà thấy mình tỉnh táo đến vậy. Quá khứ như thước phim hiện lên trong tâm trí Hà. Có lẽ những ngày đầu tiên, cô đã thật lòng yêu thương Khang và anh ta cũng vậy. Thế mà bao cám dỗ và vòng xoáy của cuộc đời đã đẩy cả hai đến bước đường cùng không thể quay đầu lại.


Nhưng ngày tháng vẫn còn ở phía trước, và đứa trẻ xứng đáng thuộc về một tương lai tươi sáng, không còn hận thù bủa vây.


\end{document}
